\documentclass[11pt]{article}
\usepackage{icmlko}
\begin{document}

\title{두 해 만에 처음 재 본 것 --- 검색 시계열은 판 $-0.0283$ 짜리고, 그건 우리가 만진 어떤 손잡이보다 여섯 배 크다}
\author{월드모델 실험실}
\date{노트 490 · 2026년 8월 3일}
\maketitle

\begin{abstract}
노트 489 가 행동 시계열 수집에 ``격차 몫 $71 \sim 91\%$''를 매겼다. 그
다음 물음은 당연한 것인데 \textbf{두 해 동안 한 번도 안 물었다} ---
\emph{그 수집이 판에 얼마짜리인가.}

트렌드 $3$ 축과 위키 $3$ 축을 \textbf{학습과 예보 양쪽에서 통째로
없애고}(애초에 안 모은 셈) 짝 씨앗 $12$ 로 쟀다.

\begin{center}
\begin{tabular}{lrrr}
\hline
설정 & 판 & 판 차 & 부호 \\
\hline
챔피언 & $+0.4828$ & --- & --- \\
위키 없이 & $+0.4757$ & $-0.0071$ & $12/12$ \\
트렌드 없이 & $+0.4622$ & $-0.0206$ & $12/12$ \\
\textbf{시계열 둘 다 없이} & $\mathbf{+0.4544}$ & $\mathbf{-0.0283}$ & $\mathbf{12/12}$ \\
대조: 태그 없이 & $+0.4592$ & $-0.0236$ & $12/12$ \\
\hline
\end{tabular}
\end{center}

날짜 통제 판도 $-0.0293$ 으로 같이 내려간다(노트 446 가드 통과). 사전
등록($0.02$ 이상)은 \textbf{통과}.

\textbf{크기를 견줘 보면} --- 두 해 동안 채택한 모형 변경은 계절 축
$+0.0041$, 특기 칸 $+0.0049$, 깊이 $6\!\to\!12$ $-0.0032$ 다.
\textbf{자료 손잡이 하나가 그 전부보다 여섯 배 크다.}
\end{abstract}

\begin{figure}[t]
\centering
\includegraphics[width=\textwidth]{figs/f1.pdf}
\caption{\textbf{가} --- 수집별 판 차. \textbf{나} --- 도메인별.
\textbf{다} --- 자료 손잡이와 모형 손잡이의 크기.}
\end{figure}

\section{거의 더해진다}

트렌드 $-0.0206$ 에 위키 $-0.0071$ 을 더하면 $-0.0277$ 이고 둘 다 빼면
$-0.0283$ 이다. \textbf{두 수집이 서로 대신하지 않는다} --- 같은 ``검색
관심''을 재는 것 같은데 겹치는 몫이 $2\%$ 뿐이다. 노트 489 의 조항(절제는
수집 단위로)이 \emph{수집들 사이}에는 안 걸린다는 뜻이다.

\section{애니는 시계열이 없는 편이 낫다}

도메인별로 보면 폭이 크다 --- 도서 $-0.1519$ · 아이돌 $-0.1170$ · 게임
$-0.0833$ 부터 \textbf{애니 $+0.0328$} 까지.

\begin{quote}
\textbf{애니는 검색 시계열을 빼면 좋아진다.}
\end{quote}

노트 481 의 조항대로 폭을 적고, 부호가 갈리므로 평균을 결론 문장에 안
쓴다. 애니는 유보 $606$ 행으로 판 무게가 크고 트렌드 덮음이 $73\%$ 라
``자료가 없어서''도 아니다. \textbf{남은 것으로 적는다} --- 노트 489 가
찾은 ``속성 표지가 붙는 행은 원래 다른 종류''의 시계열판일 수 있다.

\section{그래서 어디에 쓰나}

이 실험실은 두 해 동안 \textbf{모형}을 만졌다. 판이 $0.4766 \to 0.4856
\to 0.4828$ 로 움직였고 그 각 걸음이 $0.004$ 였다. 이 노트는 \textbf{자료
하나가 $0.028$} 이라고 말한다.

\begin{quote}
\textbf{조항} --- 새 모형 후보를 재기 전에 \textbf{그 자료를 빼면 얼마인지}
먼저 잰다. 자료 손잡이의 크기를 모르면 모형 손잡이의 크기를 읽을 눈금이
없다.
\end{quote}

그리고 곧바로 따라오는 물음 --- \emph{그 수집의 덮음은 지금 몇 퍼센트인가.}
트렌드는 아홉 도메인 합계 $41.7\%$ 다. \textbf{절반 넘는 행이 이
$-0.0283$ 짜리 자료를 못 받고 있다.} 노트 491 이 그 이유를 판다.

\section*{남은 것}

\begin{itemize}
\item \textbf{애니 $+0.0328$} --- 시계열이 해로운 유일한 큰 도메인.
\item \textbf{태그 없이 $-0.0236$} --- 대조로 넣었는데 시계열과 비슷하게
크다. 속성 수집인데 왜 큰지 따로 재야 한다(노트 489 는 속성 표지에
$27\%$ 아래를 매겼다).
\item \textbf{덮음을 올리면 이 $-0.0283$ 이 얼마나 커지나} --- 지금은
$41.7\%$ 행에서만 일하는 축의 값이다.
\end{itemize}

\end{document}
